% JAC1964-20261121, /root/arithmetic_sign.  New deduction, not a cited theorem.
\subsection*{AR-21. Reducing the whole odd Chow group to single integral curves}

The effective rational representative in AR-18 can be replaced, after
clearing one denominator, by one geometrically integral curve.  All
constructions can be made over the original number field.  This reduces
the remaining sign question to a concrete test on products of smooth
connected curves; it does not prove that sign.

\begin{lemma}[AR-21a]
Let $X$ be a smooth geometrically integral projective threefold over a
number field $K$, let $L$ be an ample line bundle on $X$, and let $V$ be an
effective one-cycle with integral coefficients.  There are a
geometrically integral curve $W\subset X$ and a positive integer $b$ such
that
\[
 [W]=[V]+b\,c_1(L)^2\quad\hbox{in }\operatorname{CH}_1(X).
 \tag{AR.21a}
\]
The normalization of $W$ is a smooth geometrically integral projective
curve over $K$.  One may in addition require $W$ to contain any
prescribed $K$-rational point of $X$.
\end{lemma}

\begin{proof}
Let $Z$ be the reduced closed subscheme supporting $V$; thus we impose
only the reduced support as a base locus, even when $V$ has
multiplicities.  For $n$ sufficiently large,
$\mathcal I_Z\otimes L^n$ is globally generated, and the linear system
$H^0(X,\mathcal I_Z\otimes L^n)$ induces an immersion on $X\setminus Z$.
The latter assertion is Stacks Project, Lemma 33.47.1, tag
\texttt{0FD4}.  Lemma 33.47.3 of the same section gives smoothness of a
general divisor away from $Z$.

We can also require regularity at the generic point of every geometric
component of $Z$.  Indeed, $Z_{\overline K}$ is generically smooth, and
its ideal has a two-dimensional conormal space at each such generic
point.  Global generation of $\mathcal I_Z\otimes L^n$ shows that the
sections whose normal derivative vanishes at that generic point form
a proper linear subspace.  Avoid the finitely many such subspaces.
These conditions and the Bertini condition form a nonempty open set
defined over $K$ in the parameter space.  Since $K$ is infinite it has
a $K$-point.  Let $S\in|L^n|$ be the resulting divisor.

Here is an explicit justification of geometric integrality.  The
hypersurface $S_{\overline K}$ is Cohen--Macaulay and is regular outside
a finite set: it is smooth off $Z_{\overline K}$ and at the generic
points of its finitely many curve components.  Serre's normality
criterion therefore makes $S_{\overline K}$ normal.  Increase $n$ in
advance so that $H^1(X_{\overline K},L^{-n})=0$.  This follows from
Serre duality and Serre vanishing applied to
$H^2(X_{\overline K},\omega_X\otimes L^n)$.  The sequence
\[
 0\longrightarrow L^{-n}\longrightarrow\mathcal O_X
 \longrightarrow\mathcal O_S\longrightarrow0
\]
then gives $H^0(S_{\overline K},\mathcal O_S)=\overline K$.
Thus $S_{\overline K}$ is connected.  The irreducible components of a
normal Noetherian scheme are disjoint, so $S_{\overline K}$ is integral.

Choose a resolution $\pi:\widetilde S\to S$ obtained by a finite
sequence of normalized blowups at closed points over the singular
points.  Such a resolution exists over $K$ by the surface resolution
theorem, Stacks Project, Theorem 54.14.5, tag \texttt{0BGJ}; its
condition (4) holds for our excellent normal surface.  The
construction in its proof is supported over the finite singular
locus.  The finite normalizations and blowups are projective, so
$\widetilde S$ is projective.  It is regular, hence smooth over the
perfect field $K$, and it is geometrically integral because it is
birational to the geometrically integral surface $S$.

There is an effective exceptional Cartier divisor $E$ on
$\widetilde S$ for which $\mathcal O_{\widetilde S}(-E)$ is
$\pi$-ample.  To see this without assuming an arbitrary exceptional
divisor works, each blowup has relatively ample
$\mathcal O(1)=\mathcal O(-E_i)$, with $E_i$ an effective exceptional
Cartier divisor (Stacks Project, Section 31.33 and Lemma 31.33.4, tag
\texttt{01OF}).  Pullback by the finite normalization preserves
relative ampleness.  For a composition, a sufficiently large positive
multiple of the pullback of the previous exceptional divisor, plus
the new exceptional divisor, works by Lemma 29.38.8, tag
\texttt{01VG}.  Induction proves the claim.  If $S$ is already smooth,
take $E=0$ and $\pi$ the identity.

By Lemma 29.38.7 of the same section, for a sufficiently large positive
integer $k$ the line bundle
\[
 H=k\pi^*(L|_S)-E
\]
is ample on $\widetilde S$.  Let $D$ be the strict transform of $V$,
with its given multiplicities.  Since $\widetilde S$ is a smooth
surface, $D$ is an effective Cartier divisor.  For $a$ sufficiently
large, $D+aH$ is very ample.  A general member
$\widetilde W\in|D+aH|$ is smooth by Bertini.  We may also take it
geometrically connected: Serre duality and vanishing give
$H^1(\widetilde S_{\overline K},\mathcal O(-D-aH))=0$ for $a\gg0$,
and the preceding exact-sequence argument applies again.  Hence
$\widetilde W$ is geometrically integral.  All these nonempty open
conditions admit a $K$-rational parameter.  We also avoid the finitely
many exceptional components.

Put $W=\pi(\widetilde W)\subset S\subset X$.  The map
$\widetilde W\to W$ is birational and proper with finite fibres, so it
is the normalization of the geometrically integral curve $W$.  Let
$f:\widetilde S\to X$ be the composite map.  The projection formula
and the fact that every exceptional curve maps to a point give
\[
 f_*D=V,\qquad f_*E=0,\qquad
 f_*\bigl(\pi^*c_1(L|_S)\bigr)=c_1(L)\cdot[S]
 =n c_1(L)^2.
\]
Proper pushforward respects rational equivalence.  Therefore
\[
 [W]=f_*[\widetilde W]=[V]+ak n c_1(L)^2
\]
in the integral Chow group, proving the assertion with $b=akn>0$.

For the last assertion, add the prescribed point $x$ to the reduced
base locus when choosing $S$.  This changes neither generic
regularity along its curve components nor smoothness away from the
base locus; a possible singularity at $x$ is allowed in the finite
singular set.  The same connectedness and normality argument
applies.  Choose a closed point $\widetilde x\in\pi^{-1}(x)$.
In choosing $\widetilde W$, use the sections of $D+aH$ vanishing at
$\widetilde x$.  For $a\gg0$ this system immerses the complement
of that finite base locus, and its first normal derivatives at the
finitely many geometric points over $\widetilde x$ are generated.
Bertini and avoidance of the corresponding proper linear
subspaces therefore give a member smooth everywhere.  Its
geometric connectedness still follows from the same vanishing for
the underlying line bundle $D+aH$.  The conditions are defined over
$K$ and have a $K$-rational parameter.  It contains
$\widetilde x$, and its image $W$ contains $x$, with exactly the
same rational-equivalence formula.  The residue field of
$\widetilde x$ is a finite extension of $K$.
\end{proof}

\begin{proposition}[AR-21b]
Let $A/K$ be an abelian threefold with a symmetric ample line bundle
$L$, and write
$\operatorname{Ce}(Z)=(Z-[-1]_*Z)/2$.
For every $z\in\operatorname{CH}^2(A)_{\mathbf Q}$ satisfying
$[-1]_*z=-z$, there exist a positive integer $M$ and a geometrically
integral curve $W\subset A$, both over $K$, such that
\[
 \operatorname{Ce}(W)=Mz
 \quad\hbox{in }\operatorname{CH}^2(A)_{\mathbf Q}.
 \tag{AR.21b}
\]
If $Lz=0$, then $L\operatorname{Ce}(W)=0$.  In the normalization of
AR-18 the exact height identity is
\[
 h_K(z,z)=\frac{H_-(W)-H_+(W)}{12h^0(A,L)M^2}.
 \tag{AR.21c}
\]
Thus semipositivity on the entire primitive grade-one group, or on its
complementary subspace $\operatorname{CH}^2(N)$, is equivalent to the
comparison $H_-(W)\ge H_+(W)$ for all geometrically integral curves
whose Ceresa class belongs to the indicated subspace.  Strict
positivity is equivalent to strict comparison for the nonzero Ceresa
classes in that subspace.
\end{proposition}

\begin{proof}
Represent $z$ by a rational cycle $Z=Z_+-Z_-$ with $Z_+,Z_-$ effective.
The effective rational cycle $V=Z_++[-1]_*Z_-$ satisfies
$\operatorname{Ce}(V)=z$ in rational Chow groups.  Choose $M>0$ so
that $MV$ has integral coefficients, and apply AR-21a to $MV$.
Symmetry of $L$ gives $[-1]_*c_1(L)^2=c_1(L)^2$, so its odd part is
zero.  Taking odd parts of (AR.21a) therefore proves (AR.21b).
Primitivity follows by multiplication by $L$.  AR-18 applied to the
actual curve $W$ gives (AR.21c), since the height pairing is quadratic
and $h_K(Mz,Mz)=M^2h_K(z,z)$.  No rational equivalence invariance of
the separate heights $H_+$ and $H_-$ is needed or asserted.  Their
difference is identified with the height only through AR-18 for this
representative.  The converse implications are immediate by applying
the proposed sign to $\operatorname{Ce}(W)$ itself.
\end{proof}

\paragraph{Field and source scope.}
No extension of the number field is required in this construction.
Geometric components were used only to specify open conditions; the
conditions are Galois invariant and their parameter spaces have
$K$-points.  The positive integer $M$ accounts for the rational cycle
coefficients and occurs quadratically in (AR.21c); no field-degree
factor is introduced.  The surface resolution and ampleness facts
are the directly read Stacks Project statements cited above.  Serre
duality, Serre vanishing, the normality criterion, and proper Chow
pushforward are used in their usual forms.  This result supplies a
single-curve test for the unresolved comparison, not a claim that
all such Ceresa classes are multiples of the Fermat class.
